\subsection{The complete full-nine-torsion range}

\begin{theorem}[B-RD23: the four remaining full-nine cases exist]
\label{thm:B-RD23}
\status{Proved; independent parent and obstructions/crystals proof/source audits passed.}
For $p=3$, an additive elliptic curve with full rational $9$-torsion
exists over a complete discretely valued field of absolute
ramification index $e$ and perfect residue field if and only if
\[
                              6\mid e,\qquad e\geq18.
                                                        \tag{B-RD23.1}
\]
The examples can all be taken over finite extensions of $\mathbf Q_3$.
In particular all four indices $24,30,42,48$ left by
Corollary~\ref{cor:B-RD21} occur, with semistability defect twelve.
\end{theorem}
\begin{proof}
Necessity follows from the cyclotomic condition and
Proposition~\ref{prop:B-RD20}. We first make the example at $e=18$
explicit over a finite local field. Let $K_0=\mathbf Q_3(i)$,
$i^2=-1$, and $E_0:y^2=x^3-x$. This has good supersingular
reduction and the endomorphism
$[i]:(x,y)\mapsto(-x,iy)$, defined over $K_0$. Put
$L=K_0(E_0[9])$, the finite Galois division field.

Every primitive nine-torsion parameter $t_0=-x/y$ has
$v_3(t_0)=1/72$. Indeed, the unramified height-two formal group's
$[3]$-series has linear coefficient three, its first unit
coefficient is in degree nine, and every intervening coefficient
is divisible by three. A nonzero three-torsion parameter therefore
has valuation $1/8$. A preimage under $[3]$ must have smaller
valuation, on which the degree-nine term has strictly least
valuation; hence its valuation is $1/72$.

On the other hand, $E_0[9]$ is free of rank one over
$\mathbf Z_3[i]/9$. To see this directly, reduction modulo three
has dimension two over $\mathbf F_3$ and the operator $i$ has
irreducible polynomial $X^2+1$; it is therefore one-dimensional
over $\mathbf F_9$. Nakayama's lemma and the equal cardinalities
then give the rank-one assertion modulo nine. Galois commutes
with $i$, so its faithful image has order at most
$|(\mathbf Z_3[i]/9)^\times|=72$. The preceding valuation
denominator forces ramification degree at least $72$.
Consequently $L/K_0$ is totally ramified of degree $72$ and this
Galois image is the full unit group.

Let $\sigma\in\operatorname{Gal}(L/K_0)$ act on $E_0[9]$ by $[i]$,
and let $K=L^{\langle\sigma\rangle}$. Choose a primitive point
$Q\in E_0[9]$ and put $u=-x(Q)/y(Q)$. It is a uniformizer of
$L$, and the exact coordinate formula for $[i]$ gives
$\sigma u=iu$. Thus $a=u^4\in K$ has $v_K(a)=1$. The curve
\[
                         E:y^2=x^3-ax
\]
has full rational nine-torsion: its points obtained from
$P\in E_0[9]$ are $(u^2x(P),u^3y(P))$, and both coordinates are
fixed by $\sigma$. The absolute ramification index of $K$ is
$72/4=18$. The curve becomes $E_0$ over $L$, while
$v_K(\Delta(E))=v_K(64a^3)=3$ forces exact defect four. This gives
the promised finite local-field example at the endpoint.

We next give a uniform explicit construction for every multiple
$e=6d>18$.

Choose $i^2=-1$, a primitive ninth root of unity $\eta$, and put
$\zeta=\eta^3$. Let
\[
 K=\mathbf Q_3(i,\eta,\pi),\qquad \pi^d=\eta-1.
                                                        \tag{B-RD23.2}
\]
The field $\mathbf Q_3(i)$ is the unramified quadratic extension,
and $\eta-1$ is a uniformizer of its degree-six cyclotomic
extension. The final polynomial in (B-RD23.2) is Eisenstein.
Thus $K$ has absolute ramification $e$, residue field $\mathbf F_9$,
and uniformizer $\pi$. Write $v(\pi)=1$.

Put
\[
 h=1+\pi^2,\quad D=3(4h^3-1),\quad
 y=\frac{-3+\sqrt D}{2(1-h^3)}.                         \tag{B-RD23.3}
\]
Here $v(h^3-1)=6$, $v(4h^3-1)=6$, and $v(D)=e+6$.
Moreover $3/\pi^e\equiv-1\pmod\pi$, by the Eisenstein
cyclotomic equation for $\eta-1$. Consequently
$D/\pi^{e+6}\equiv-1\pmod\pi$. This unit has a square root in
$K$ lifting either $i$ or $-i$, by ordinary Hensel lifting, since
two is a unit. This specifies the square root in (B-RD23.3).
The quadratic formula gives
\[
       h^3=1+\frac3y+\frac3{y^2},\qquad
       r:=v(y)=\frac e2-3.                             \tag{B-RD23.4}
\]
In particular $r>0$. Define
\[
 \lambda=-\eta^2,\qquad c=\frac{\lambda}{\zeta-1},\qquad
 d_0=(1+y)^3-1=y^3h^3,\qquad
 t=1+\frac{1-\zeta}{d_0}.                             \tag{B-RD23.5}
\]
Since $(\zeta-1)^2=-3\zeta$, one has
$c^3=(1-\zeta)/9$. Therefore
\[
 \frac{t-1}{9}=\left(\frac{c}{yh}\right)^3,
 \qquad
 \frac{t-\zeta}{9}
   =\left(\frac{c(1+y)}{yh}\right)^3.                 \tag{B-RD23.6}
\]

We show that $(t-\zeta^2)/9$ is also a cube. Its ratio with the
first expression in (B-RD23.6) is
\[
              a=1-\zeta^2d_0=1+\lambda^3d_0.
\]
Use $\beta_0=1+\lambda y$ as a first approximation to a cube root.
Then
\[
 a-\beta_0^3
      =3(\lambda^3-\lambda)y
                   +3(\lambda^3-\lambda^2)y^2.
                                                        \tag{B-RD23.7}
\]
The first coefficient in parentheses has valuation $e/6$ and the
second is a unit. Since $r>e/6$, the valuation in (B-RD23.7) is
\[
             e+q,\qquad q=e/6+r=2e/3-3>e/2.
\]
The polynomial
\[
       \frac{(\beta_0+\pi^q Z)^3-a}{3\pi^q}
\]
has integral coefficients. Its linear coefficient is the unit
$\beta_0^2$, its quadratic coefficient is $\beta_0\pi^q$, and
its cubic coefficient has valuation $2q-e>0$. Its reduction is
therefore a degree-one polynomial, which has a simple root over
$\mathbf F_9$. Hensel lifting gives $\beta\in K$ with
$\beta^3=a$. This argument uses the displayed rescaling; it does
not invoke the generally inapplicable stronger inequality
$v(a-\beta_0^3)>2v(3)$.

Now use the original full-level modular-curve equations of
Fisher \cite[Lemma~4.2(i), p.8; Theorem~4.3, equation~(9), p.9]
{Fisher2015Nine}. They state that the family over $Y(3)$ is
\[
             E_t:\quad y_0^2+3tx_0y_0+(t^3-1)y_0=x_0^3,
                                                        \tag{B-RD23.8}
\]
and the cover $X(9)\to X(3)$ is given by the three equations
\[
     t-\theta=9(u+v\theta+w\theta^2)^3,
                     \qquad \theta=1,\zeta,\zeta^2.    \tag{B-RD23.9}
\]
The three cube roots supplied above are respectively
$c/(yh)$, $c(1+y)/(yh)$, and $c\beta/(yh)$. An invertible Fourier
matrix over $K$ gives $u,v,w$ with precisely these three values.
Thus (B-RD23.9) gives a $K$-point of $Y(9)$. By the moduli
interpretation in Section~2, p.4 of that source, it gives
$E_t[9]\simeq\mu_9\oplus\mathbf Z/9$. Since $\mu_9\subset K$,
the entire $9$-torsion is rational. The parameter is not a cusp,
because the valuation calculation below gives $v(t)<0$.

Finally, (B-RD23.4)--(B-RD23.5) give $v(t)=9-e<0$. The exact
invariants of (B-RD23.8) are
\[
 c_4=9t(t^3+8),\qquad \Delta=27(t^3-1)^3,\qquad
 j=\frac{27t^3(t^3+8)^3}{(t^3-1)^3}.
\]
Hence
\[
                   v(j)=27,\qquad v(\Delta)=81-6e
                                             \equiv9\pmod {12}.
\]
The integral $j$ gives potential good reduction
\cite[Section~2, p.496]{SerreTate1968Good}, and the discriminant
congruence excludes good reduction. Thus the reduction is additive.
An explicit integral equation is obtained by the substitution
$x_0=\pi^{2(9-e)}X$, $y_0=\pi^{3(9-e)}Y$: its two nonzero
coefficients have valuations $e$ and zero, and its discriminant
has valuation $6e-27$.

At any of $e=24,30,42,48$, the discriminant congruence forces the
minimal good inertia order $m$ to be divisible by four. Its possible
orders in characteristic three are therefore four and twelve.
The sharp tame full-nine criterion excludes four at these indices,
so $m=12$. All displayed finite polynomial identities, the Fourier
inversion, and the valuation inequalities at the four previously
undecided indices are checked in
\path{code/verify_full9_construction.py}. The square-root and cubic
Hensel steps are justified above, rather than treated as numerical
approximations.
\end{proof}

\begin{remark}[Passage to all higher levels]
This first settles level $n=2$ of the wild higher-torsion question,
including every exceptional band with $k=1$ in Corollary~\ref{cor:B-RD21}.
The first higher case beyond this construction is full $27$-torsion
at $e=216$. It is subsequently realized by the wild trace-invariant
construction of Theorem~\ref{thm:B-RD26-wild}, which settles every
higher wild level and yields Corollary~\ref{cor:B-RD27-exact-three}.
\end{remark}
